令 \(A=2/p_c\) 和 \(L(R)=\log(2+R)\)。\ref{ms:cor-9-1} 中指数为 \(2A\) 的 Kakeya 估计，给出了界 \(C L(1/\rho)^{2A}\)，见 \ref{ms:nikodym-distal}，针对的是 \(D_\rho\)，即在距离介于 \(1/2\) 与 \(1\) 之间的射线段邻域上取上确界，距离从 \(x\) 起算。

任取一条经过 \(x\) 的单位线段。线段上距中心不超过 \(2\delta\) 的点（中心为 \(x\)），连同这些点的 \(\delta\) 邻域，均落在 \(B(x,3\delta)\) 中。沿两个射线方向，剩余部分分别落在由距离范围 \([s/2,s]\) 确定的若干区间的并集中，其中 \(s=2^{-j}\)、\(j\ge0\)、\(s\ge\delta\)。对这样的 \(s\)，射线区间的 \(\delta\) 邻域体积为 \(\lesssim_n s\delta^{n-1}\)，因为 \(\delta\le s\)。整个管的体积与 \(\delta^{n-1}\) 可比。因此，球和射线区间的贡献，各自至多为 \(\delta\) 和 \(s\) 的常数倍再乘以各自的极大平均值。由 Young 不等式，球平均的 \(L^n\) 算子范数为一。尺度 \(s\) 上的射线极大平均经伸缩共轭于 \(D_{\delta/s}\)，故其范数由 \(C L(s/\delta)^{2A}\) 控制。因此
\[
 \|N_\delta\|_{L^n\to L^n}
 \lesssim \delta+\sum_{\substack{s=2^{-j},\ j\ge0\\ s\ge\delta}}sL(s/\delta)^{2A}
 \lesssim L(1/\delta)^{2A}.
\]
